\subsection*{先检测，再写回}
前馈层可以看成两阶段：第一组权重对输入做许多测量，非线性决定这些测量怎样激活，第二组权重把激活写回表示空间。一个简化单元先计算 $k^Tx$，再沿方向 $v$ 输出相应强度，贡献为 $v(k^Tx)$。

取 $k=(1,2),v=(3,-1),x=(2,1)$，检测值为四，写回向量是 $(12,-4)$。先看检测方向，再看写回方向，可以避免把同一神经元的输入权重和输出权重误认成同一种角色。

\subsection*{外积显示单个通道能做多大的事}
把表达式重新加括号，是 $(vk^T)x$。外积矩阵的每列都与 $v$ 平行，所以秩至多一：单个线性通道只能沿一条输出方向写信息。多个通道相加能形成更丰富的作用，而非线性又使不同输入激活不同组合。

在 Transformer 中，前馈层通常逐位置处理表示，位置之间的通信主要由注意力承担；不过每个位置的表示此前已经包含上下文，所以“逐位置”不等于“完全不知上下文”。

把通道解释成某个特征或关联的检测器，是有用的构造性图景，却不是对真实模型的唯一解读。特征可能分散在多个单元，也可能在一个单元中叠加。练习时先精确核对外积的计算，再把“网络可以实现什么”与“训练后实际实现什么”分开。
